\documentclass[11pt]{article}

\usepackage[T1]{fontenc}
\usepackage{amsmath,amssymb,amsthm}
\usepackage[margin=1in]{geometry}
\usepackage[colorlinks=true,linkcolor=blue,urlcolor=blue]{hyperref}

\newtheorem{theorem}{Theorem}
\newtheorem{proposition}[theorem]{Proposition}
\newtheorem{corollary}[theorem]{Corollary}
\newtheorem{conjecture}[theorem]{Conjecture}
\theoremstyle{remark}
\newtheorem{remark}[theorem]{Remark}

\newcommand{\R}{\mathbb R}
\newcommand{\Pp}{\mathbb P}
\newcommand{\norm}[1]{\left\|#1\right\|}
\newcommand{\ip}[2]{\left\langle #1,#2\right\rangle}
\newcommand{\weakto}{\rightharpoonup}
\newcommand{\Var}{\operatorname{Var}}
\setlength{\emergencystretch}{2em}

\title{Diagonal Gaussian Flux and Exact Recycling\\
at the q4 Nonlinear Entrance}
\author{John Lawson and Codex}
\date{25 July 2026}

\begin{document}
\maketitle

\begin{abstract}
This report continues the conditional energy-efficient exact q4 branch
of the repository.  Its reversed full-defect component \(v\) has weak
boundary trace zero but limiting energy \(d>0\), and all of that energy
is regenerated by a projected quadratic Duhamel term.  For
\[
 F:=\Pp((U\cdot\nabla)v),\qquad S(r):=e^{\nu r\Delta},
\]
the moving heat field \(S(\tau-s)v(s)\) evolves by
\(-S(\tau-s)F(s)\): its viscosity terms cancel exactly.  Consequently
half of the entrance energy is the current-time diagonal Gaussian flux
\[
 \frac12\norm{v(\tau)}_2^2
 =
 \int_0^\tau-\ip{F(s)}{S(2(\tau-s))v(s)}\,ds.
\]
We prove the exact recycling identity between nested terminal
detectors.  Pairwise disjoint heat-age bands also satisfy a genuine
\(\ell^1\) estimate, but its right side is
\(M\norm{\nabla v}_2^2\), whose time integral is already forced to
diverge in q4.

A coherent positive scalar kernel satisfies the endpoint, diagonal
energy, recycling, magnitude, variation, and power--log constraints
while concentrating arbitrarily much work in history shared by nested
detectors.  Its constant-energy core has an exact signed Laguerre
heat-spectrum representation.  The kernel is not a velocity field or a
Navier--Stokes triad.  Thus abstract Gaussian-band orthogonality does
not close q4; a new theorem must exploit the actual quadratic spectral
relation or produce a finite same-trajectory charge.  All findings are
conditional repository deductions pending external review.  No Clay
alternative is proved.
\end{abstract}

\section{Conditional setting}

Let \(v,c\) be solenoidal fields on
\((0,\tau_0)\times\R^3\), smooth at every positive time, and put
\[
 U:=v+c,\qquad F:=\Pp((U\cdot\nabla)v).
\]
Assume
\begin{equation}
\label{eq:evolution}
 \partial_s v-\nu\Delta v+F=0
\end{equation}
and the entrance laws
\begin{align}
\label{eq:weak}
 v(s)&\weakto0
 &&\text{in \(L^2_\sigma(\R^3)\) as \(s\downarrow0\)},\\
\label{eq:energy}
 \norm{v(s)}_2^2+
 2\nu\int_0^s\norm{\nabla v(a)}_2^2\,da&=d>0.
\end{align}
The preceding boundary mild theorem supplies the stronger fixed-terminal
statement
\begin{equation}
\label{eq:heat-boundary}
 S(\tau-s)v(s)\longrightarrow0
 \quad\text{strongly in \(L^2\) as \(s\downarrow0\)}
\end{equation}
for every \(\tau\in(0,\tau_0)\), where
\[
 S(r):=e^{\nu r\Delta}.
\]
Write
\begin{equation}
\label{eq:MY}
 M(s):=\norm{U(s)}_{L^{3,\infty}},
 \qquad
 Y(s):=\norm{\nabla v(s)}_2.
\end{equation}
The q4 input includes finite unweighted dissipation but divergent
amplitude-weighted dissipation:
\begin{equation}
\label{eq:weighted-divergence}
 \int_0^hY(s)^2\,ds<\infty,
 \qquad
 \int_0^hM(s)Y(s)^2\,ds=\infty
 \quad(0<h<\tau_0).
\end{equation}

\begin{remark}
These assumptions are not known for arbitrary Clay-admissible data.
They arise only after the repository's conditional q4 reduction.
\end{remark}

\section{The moving heat identity}

For fixed \(\tau\), define
\begin{equation}
\label{eq:moving-field}
 w_\tau(s):=S(\tau-s)v(s),
 \qquad 0<s\leq\tau,
\end{equation}
and for heat age \(r\geq0\) define
\begin{equation}
\label{eq:flux}
 \Pi_r(s):=-\ip{F(s)}{S(2r)v(s)}.
\end{equation}

\begin{theorem}[Viscosity-free moving heat evolution]
\label{thm:moving}
For \(0<s<\tau\),
\begin{equation}
\label{eq:moving-derivative}
 \partial_s w_\tau(s)=-S(\tau-s)F(s).
\end{equation}
Moreover,
\begin{equation}
\label{eq:diagonal}
 \boxed{
 \frac12\norm{v(\tau)}_2^2
 =
 \int_0^\tau\Pi_{\tau-s}(s)\,ds
 }
\end{equation}
and
\begin{equation}
\label{eq:zero-age}
 \Pi_0(s)=0.
\end{equation}
\end{theorem}

\begin{proof}
Differentiating \eqref{eq:moving-field} and using
\eqref{eq:evolution} gives
\[
\begin{aligned}
 \partial_sw_\tau
 &=-\nu\Delta S(\tau-s)v
   +S(\tau-s)(\nu\Delta v-F)\\
 &=-S(\tau-s)F,
\end{aligned}
\]
which proves \eqref{eq:moving-derivative}.  Self-adjointness and the
semigroup law give
\[
\begin{aligned}
 \frac{d}{ds}\frac12\norm{w_\tau(s)}_2^2
 &=-\ip{S(\tau-s)F(s)}{S(\tau-s)v(s)}\\
 &=-\ip{F(s)}{S(2(\tau-s))v(s)}
 =\Pi_{\tau-s}(s).
\end{aligned}
\]
At the lower endpoint, \eqref{eq:heat-boundary} gives
\(w_\tau(s)\to0\) strongly.  At the upper endpoint,
\(w_\tau(\tau)=v(\tau)\).  Integration proves
\eqref{eq:diagonal}.

Finally, the Leray projection is self-adjoint on solenoidal test fields,
and incompressibility gives
\[
 \Pi_0(s)
 =-\int_{\R^3}(U\cdot\nabla)v\cdot v\,dx
 =-\frac12\int_{\R^3}U\cdot\nabla|v|^2\,dx
 =0.
\]
\end{proof}

\begin{remark}
The cutoff in \(w_\tau\) opens at the physical heat rate, so its
derivative cancels viscosity.  Formula \eqref{eq:diagonal} contains no
future-time test field and no pressure term.
\end{remark}

\begin{proposition}[Continuous heat-age bands]
\label{prop:bands}
For every \(r>0\),
\begin{equation}
\label{eq:continuous-band}
 \boxed{
 \Pi_r(s)
 =
 \ip{F(s)}{(I-S(2r))v(s)}
 =
 -2\nu\int_0^r
 \ip{F(s)}{\Delta S(2\rho)v(s)}\,d\rho .
 }
\end{equation}
\end{proposition}

\begin{proof}
Equation \eqref{eq:zero-age} and the definition of \(\Pi_r\) yield the
first equality.  Since
\[
 I-S(2r)=-2\nu\int_0^r\Delta S(2\rho)\,d\rho,
\]
the second follows.
\end{proof}

\begin{proposition}[Pointwise heat estimate]
\label{prop:pointwise}
For \(r>0\),
\begin{equation}
\label{eq:pointwise}
 |\Pi_r(s)|
 \lesssim
 \sqrt d\,(\nu r)^{-1/2}M(s)Y(s).
\end{equation}
\end{proposition}

\begin{proof}
Lorentz Holder, boundedness of \(\Pp\), and heat smoothing give
\[
 \norm{F(s)}_{L^{6/5,2}}\lesssim M(s)Y(s),
 \qquad
 \norm{S(2r)v(s)}_{L^{6,2}}
 \lesssim(\nu r)^{-1/2}\norm{v(s)}_2.
\]
The energy identity bounds \(\norm{v(s)}_2\) by \(\sqrt d\).
Lorentz duality now proves \eqref{eq:pointwise}.
\end{proof}

\section{Exact recycling}

\begin{theorem}[Nested terminal recycling]
\label{thm:recycling}
If \(0<\sigma<\tau<\tau_0\), then
\begin{equation}
\label{eq:recycling}
 \boxed{
 \begin{aligned}
 &\int_0^\sigma
 \bigl(\Pi_{\sigma-s}(s)-\Pi_{\tau-s}(s)\bigr)\,ds\\
 &\qquad=
 \int_\sigma^\tau\Pi_{\tau-s}(s)\,ds
 +\nu\int_\sigma^\tau Y(s)^2\,ds.
 \end{aligned}
 }
\end{equation}
\end{theorem}

\begin{proof}
Set
\[
 \Phi(t):=\int_0^t\Pi_{t-s}(s)\,ds
 =\frac12\norm{v(t)}_2^2.
\]
Splitting the difference at \(\sigma\) gives
\[
\begin{aligned}
 \Phi(\tau)-\Phi(\sigma)
 &=
 \int_0^\sigma
 \bigl(\Pi_{\tau-s}(s)-\Pi_{\sigma-s}(s)\bigr)\,ds\\
 &\quad+\int_\sigma^\tau\Pi_{\tau-s}(s)\,ds.
\end{aligned}
\]
On the other hand, \eqref{eq:energy} gives
\[
 \Phi(\tau)-\Phi(\sigma)
 =-\nu\int_\sigma^\tau Y(s)^2\,ds.
\]
Rearrangement proves \eqref{eq:recycling}.
\end{proof}

\begin{remark}
For \(s<\sigma\), the left side of \eqref{eq:recycling} is the
nonlinear pairing through the Gaussian band between heat ages
\(\sigma-s\) and \(\tau-s\).  The identity records exactly how shared
early history pays for later-block flux and viscous loss.
\end{remark}

\section{Orthogonality and its exact obstruction}

\begin{theorem}[Disjoint Gaussian-band budget]
\label{thm:l1}
Fix \(s>0\).  Let \(\{[a_k,b_k]\}\) be finitely many pairwise disjoint
subintervals of \([0,\infty)\), with \(a_k<b_k\), and set
\[
 B_k:=S(2a_k)-S(2b_k).
\]
Then
\begin{equation}
\label{eq:l1}
 \boxed{
 \sum_k\left|\ip{F(s)}{B_kv(s)}\right|
 \lesssim M(s)Y(s)^2.
 }
\end{equation}
\end{theorem}

\begin{proof}
The Fourier symbol
\[
 \beta_k(\xi)
 =
 e^{-2\nu a_k|\xi|^2}
 -e^{-2\nu b_k|\xi|^2}
\]
is nonnegative.  Since the heat-age intervals are disjoint,
\[
 \sum_k\beta_k(\xi)\leq1
 \qquad(\xi\in\R^3).
\]
Indeed, each difference is the integral of the positive derivative
density \(2\nu|\xi|^2e^{-2\nu r|\xi|^2}\) over its heat-age interval.

Let \(C_k:=B_k^{1/2}\).  Cauchy--Schwarz in \(k\), followed by
Plancherel, yields
\[
\begin{aligned}
 \sum_k|\ip{F}{B_kv}|
 &\leq
 \left(\sum_k\norm{C_kF}_{\dot H^{-1}}^2\right)^{1/2}
 \left(\sum_k\norm{C_kv}_{\dot H^1}^2\right)^{1/2}\\
 &\leq\norm{F}_{\dot H^{-1}}\norm{\nabla v}_2.
\end{aligned}
\]
For a solenoidal \(\phi\in\dot H^1\), Lorentz Holder and the refined
Sobolev embedding \(\dot H^1\hookrightarrow L^{6,2}\) give
\[
\begin{aligned}
 |\ip{F}{\phi}|
 &=\left|\int_{\R^3}U_jv_i\partial_j\phi_i\,dx\right|\\
 &\lesssim
 \norm{U}_{L^{3,\infty}}\norm{v}_{L^{6,2}}
 \norm{\nabla\phi}_2\\
 &\lesssim M Y\norm{\nabla\phi}_2.
\end{aligned}
\]
Since \(F\) is solenoidal, testing against the Leray projection handles
general \(\phi\).  Hence
\(\norm{F}_{\dot H^{-1}}\lesssim MY\), proving
\eqref{eq:l1}.
\end{proof}

\begin{corollary}[Finite-range recycling variation]
\label{cor:variation}
Let
\(\tau_0>\tau_1>\cdots>0\), and define
\[
 \mathcal R_j:=
 \int_0^{\tau_{j+1}}
 \bigl(\Pi_{\tau_{j+1}-s}(s)
       -\Pi_{\tau_j-s}(s)\bigr)\,ds.
\]
For finite \(J\leq K\),
\begin{equation}
\label{eq:recycling-variation}
 \boxed{
 \sum_{j=J}^K|\mathcal R_j|
 \lesssim
 \int_0^{\tau_{J+1}}M(s)Y(s)^2\,ds.
 }
\end{equation}
\end{corollary}

\begin{proof}
For each fixed \(s\), the intervals
\[
 [\tau_{j+1}-s,\tau_j-s],
 \qquad s<\tau_{j+1},
\]
are pairwise disjoint.  Apply Theorem~\ref{thm:l1}, integrate in \(s\),
and use the triangle inequality.
\end{proof}

\begin{remark}[The obstruction]
Corollary~\ref{cor:variation} is the desired \(\ell^1\) non-reuse
estimate for Gaussian bands.  It does not close q4 because
\eqref{eq:weighted-divergence} makes its natural budget infinite on
every entrance interval.  The estimate has recovered exactly the
already-divergent amplitude-weighted dissipation.
\end{remark}

\section{A coherent scalar survivor}

Let
\[
 \ell(s):=\log(e/s),\qquad
 M_\sharp(s):=s^{-2/11}\ell(s)^{-2/11},\qquad
 Y_\sharp(s):=s^{-9/22}\ell(s)^{1/11}.
\]
Then, as \(t\downarrow0\),
\begin{align}
\label{eq:sharp-tails}
 \int_0^tM_\sharp(s)^2\,ds
 &\asymp t^{7/11}\ell(t)^{-4/11},\\
 \int_0^tY_\sharp(s)^2\,ds
 &\asymp t^{2/11}\ell(t)^{2/11},\\
\label{eq:sharp-action}
 M_\sharp(s)Y_\sharp(s)^2&=s^{-1}.
\end{align}
Put
\[
 D_\sharp(t):=\int_0^tY_\sharp(s)^2\,ds,
 \qquad
 E_\sharp(t):=d-2\nu D_\sharp(t),
\]
and restrict to a terminal interval on which \(E_\sharp>0\).
For an integer \(n\geq1\), define
\begin{equation}
\label{eq:scalar-kernel}
 \Pi_n^\sharp(r,s):=
 \frac{(n+1)E_\sharp(r+s)}2
 \frac{r^n}{(r+s)^{n+1}}.
\end{equation}

\begin{proposition}[Coherent diagonal reuse]
\label{prop:survivor}
The kernel \eqref{eq:scalar-kernel} is positive for \(r>0\), vanishes
at \(r=0\), and satisfies
\begin{equation}
\label{eq:scalar-diagonal}
 \int_0^\tau\Pi_n^\sharp(\tau-s,s)\,ds
 =\frac12E_\sharp(\tau).
\end{equation}
It therefore satisfies the scalar version of
Theorem~\ref{thm:recycling}.  Near the entrance,
\begin{equation}
\label{eq:scalar-magnitude}
 \Pi_n^\sharp(r,s)
 \lesssim_n r^{-1/2}M_\sharp(s)Y_\sharp(s),
\end{equation}
and
\begin{equation}
\label{eq:scalar-variation}
 \Var_{r>0}\Pi_n^\sharp(r,s)
 \lesssim_n s^{-1}
 =M_\sharp(s)Y_\sharp(s)^2.
\end{equation}
For \(0<q<1\),
\begin{equation}
\label{eq:reuse-fraction}
 \boxed{
 \frac{\displaystyle
  \int_0^{q\tau}\Pi_n^\sharp(\tau-s,s)\,ds}
 {\displaystyle
  \int_0^\tau\Pi_n^\sharp(\tau-s,s)\,ds}
 =1-(1-q)^{n+1}.
 }
\end{equation}
\end{proposition}

\begin{proof}
Positivity and vanishing at \(r=0\) are immediate.  Along
\(r+s=\tau\),
\[
\begin{aligned}
 \int_0^\tau\Pi_n^\sharp(\tau-s,s)\,ds
 &=
 \frac{(n+1)E_\sharp(\tau)}{2\tau^{n+1}}
 \int_0^\tau(\tau-s)^n\,ds\\
 &=\frac12E_\sharp(\tau),
\end{aligned}
\]
which proves \eqref{eq:scalar-diagonal}.  Since
\[
 E_\sharp(\tau)-E_\sharp(\sigma)
 =-2\nu\int_\sigma^\tau Y_\sharp(s)^2\,ds,
\]
the same subtraction used in Theorem~\ref{thm:recycling} proves the
recycling identity.

For \eqref{eq:scalar-magnitude}, set \(x=r/s\).  Up to a constant
depending on \(n,d,\nu\), the ratio of the two sides is bounded by
\[
 s^{1/11}\ell(s)^{1/11}
 \frac{x^{n+1/2}}{(1+x)^{n+1}}.
\]
The \(x\)-factor is bounded, and the \(s\)-factor tends to zero.

For fixed \(s\), the core is
\[
 \frac{r^n}{(r+s)^{n+1}}
 =s^{-1}\frac{(r/s)^n}{(1+r/s)^{n+1}}.
\]
It is unimodal and has total \(r\)-variation \(C_n/s\).
The positive bounded factor \(E_\sharp(r+s)\) is monotone and has
bounded total variation.  The product rule for bounded variation proves
\eqref{eq:scalar-variation}.

Finally, the diagonal factor \(E_\sharp(\tau)\) cancels in the quotient,
and direct integration gives
\[
 \frac{n+1}{\tau^{n+1}}\int_0^{q\tau}(\tau-s)^n\,ds
 =1-(1-q)^{n+1}.
\]
\end{proof}

\begin{remark}
For fixed \(q\), the fraction \eqref{eq:reuse-fraction} tends to one as
\(n\to\infty\).  Thus the scalar constraints permit almost all diagonal
work to lie in the early interval shared by two nested detectors.
\end{remark}

\begin{proposition}[Laguerre heat spectrum of the core]
\label{prop:laguerre}
Let
\[
 L_n(z):=\sum_{k=0}^n\binom nk\frac{(-z)^k}{k!}.
\]
Then
\begin{equation}
\label{eq:laguerre}
 \boxed{
 \frac{r^n}{(r+s)^{n+1}}
 =
 \int_0^\infty
 e^{-r\lambda}e^{-s\lambda}L_n(s\lambda)\,d\lambda.
 }
\end{equation}
\end{proposition}

\begin{proof}
Expand \(r^n=((r+s)-s)^n\):
\[
 \frac{r^n}{(r+s)^{n+1}}
 =
 \sum_{k=0}^n\binom nk
 \frac{(-s)^k}{(r+s)^{k+1}}.
\]
The gamma integral gives
\[
 \frac1{(r+s)^{k+1}}
 =
 \frac1{k!}\int_0^\infty
 \lambda^ke^{-(r+s)\lambda}\,d\lambda.
\]
Substitution and collection of the polynomial prove
\eqref{eq:laguerre}.  At \(r=0\), the signed density integrates to zero
for \(n\geq1\), matching endpoint cancellation.
\end{proof}

\begin{remark}[Scope of the survivor]
Proposition~\ref{prop:laguerre} represents only the
constant-energy core \(r^n/(r+s)^{n+1}\).  It does not represent the
full factor \(E_\sharp(r+s)\).  More importantly,
\(\Pi_n^\sharp\) is not a spatial field, does not impose
\(F=\Pp((U\cdot\nabla)v)\), and contains no coupled pressure or
same-trajectory triad genealogy.  It is a countermodel only to closure
from the scalar identities and estimates proved above.
\end{remark}

\section{Round disposition}

\subsection*{Robust conditional findings, subject to review}

\begin{enumerate}
 \item The moving heat field evolves only by filtered nonlinearity;
 viscosity cancels exactly.
 \item Half of the entrance energy is the current-time diagonal Gaussian
 flux \eqref{eq:diagonal}.
 \item The flux vanishes at zero heat age and is an accumulated signed
 heat-band pairing.
 \item Nested terminal detectors obey the exact recycling law
 \eqref{eq:recycling}.
 \item Disjoint positive Gaussian bands have the \(\ell^1\) budget
 \eqref{eq:l1}.
 \item The corresponding time budget is \(MY^2\), which q4 already
 forces to diverge.
 \item The scalar kernel \eqref{eq:scalar-kernel} meets every scalar
 diagonal, recycling, magnitude, variation, and q4 power--log
 constraint while permitting arbitrarily strong shared-history reuse.
 \item The survivor's constant-energy core has the exact signed
 Laguerre representation \eqref{eq:laguerre}.
\end{enumerate}

\subsection*{Closed shortcut}

Positive Gaussian band multipliers and abstract frequency orthogonality
do prevent unrestricted double counting, but only against the
amplitude-weighted dissipation \(\int MY^2\).  Since that quantity is
necessarily infinite in the surviving q4 branch, this orthogonality
cannot by itself contradict a full entrance defect.

\subsection*{Open obligations}

\begin{enumerate}
 \item Write the radial heat spectral density of the actual flux and
 expose its exact quadratic Fourier-triad formula.
 \item Use incompressibility, phase, frequency genealogy, or pressure
 coupling to exclude coherent Laguerre-type reuse on one trajectory.
 \item Replace the \(MY^2\) band budget by a finite same-trajectory
 charge, or prove a fixed fresh component paid by unweighted
 dissipation.
 \item Make the strong-trace drift contribution negligible at the q4
 moving heat clock and isolate a coercive self-interaction law.
 \item Combine the Hausdorff-one-null initial energy measure with the
 moving heat evolution to obtain a finite spatial-capacity charge.
 \item Submit the conditional chain and this round to independent
 external mathematical review.
 \item Treat slower clocks, divergent normalised energy, R3B, and the
 other Clay alternatives separately.
\end{enumerate}

\begin{conjecture}[No coherent nonlinear Gaussian reuse]
\label{conj:nonreuse}
No one same-trajectory q4 entrance with finite unweighted dissipation can
have its actual flux
\[
 \Pi_r(s)
 =
 -\ip{\Pp((U\cdot\nabla)v)(s)}{S(2r)v(s)}
\]
asymptotically realise a coherent family of the form
\eqref{eq:scalar-kernel} across every record block.
\end{conjecture}

\begin{remark}
Conjecture~\ref{conj:nonreuse} is not proved.  Its proof appears to
require a genuinely new theorem about the actual quadratic
same-trajectory spectral density.  No claim of regularity, breakdown,
energy equality, q4 closure, or a Clay A--D resolution is made.
\end{remark}

\section*{Conclusion}

This round converts a future-dependent work identity into an exact
current-time Gaussian flux and proves the strongest natural abstract
non-reuse estimate for disjoint heat-age bands.  The result is genuine
progress in localisation of the obstruction, but not closure: the
orthogonality charge is exactly the divergent q4 weighted action, and a
coherent scalar survivor saturates all resulting constraints.  The next
gate is no longer generic heat calculus.  It is the actual quadratic
triad structure, or a new finite same-trajectory charge.  The Clay
problem remains unsolved.

\end{document}
